\chapter{Objetivos}
\label{cap:objetivos}

Este capítulo presenta el objetivo principal del trabajo y los objetivos específicos que se
derivan de él, así como las restricciones y condicionantes que han guiado el desarrollo del
proyecto.

% ─────────────────────────────────────────────────────────────
\section{Objetivo principal}
\label{sec:objetivo_principal}
% ─────────────────────────────────────────────────────────────

El objetivo principal de este Trabajo de Fin de Grado es diseñar, implementar y evaluar
estrategias de \textit{prompting} para analizar hasta qué punto los modelos de lenguaje de
gran escala pueden resolver puzzles de olimpiadas lingüísticas a partir de un conjunto
reducido de ejemplos, apoyándose en razonamiento metalingüístico y sin recibir información
explícita previa sobre la lengua en cuestión.

% ─────────────────────────────────────────────────────────────
\section{Objetivos específicos}
\label{sec:objetivos_especificos}
% ─────────────────────────────────────────────────────────────

De este objetivo principal se derivan los siguientes objetivos específicos:

\begin{enumerate}
	
	\item \textbf{Revisar el estado del arte} en razonamiento lingüístico con modelos de
	lenguaje de gran escala, con especial atención a los benchmarks derivados de la
	Olimpiada Internacional de Lingüística (\textsc{Linguini}, LingOly, modeLing) y a
	las estrategias de prompting más relevantes de la literatura reciente.
	
	\item \textbf{Diseñar e implementar un control exploratorio previo y cinco estrategias
		de prompting} diferenciadas y comparables entre sí:
	
	\begin{itemize}
		\item Un \textit{control exploratorio previo} (M0) que permite detectar posibles
		indicios de conocimiento de la lengua antes del experimento.
		
		\item Una estrategia \textit{baseline} (M1) que sirve como punto de referencia
		directo con la literatura existente, presentando todos los ejemplos de
		entrenamiento de forma simultánea.
		
		\item Una estrategia de \textit{presentación incremental acumulativa} (M2), en la
		que se incorporan progresivamente nuevos ejemplos y se construye de forma
		explícita y revisable un diccionario y un conjunto de reglas gramaticales.
		
		\item Una estrategia \textit{paso a paso} (M3) ordenada por nivel lingüístico
		(léxico, fonología, morfosintaxis y sintaxis), inspirada en el trabajo de
		\citet{zhu2025evaluating}.
		
		\item Una estrategia de \textit{inspiración humana} (M4) que proporciona al modelo
		el razonamiento explícito de una persona resolviendo un puzzle diferente, con el
		objetivo de guiar el proceso inductivo.
		
		\item Una estrategia de \textit{autocorrección iterativa} (M5), en la que el
		modelo revisa sus propias hipótesis en un segundo turno buscando contradicciones
		internas, inspirada en el paradigma \textit{Self-Refine}
		\citep{madaan2023selfrefine}.
	\end{itemize}
	
	\item \textbf{Plantear e implementar una extensión combinatoria de la estrategia
		incremental} (M2.2), en la que se evalúan todas las combinaciones posibles de frases
	de entrenamiento para cada tamaño de subconjunto. Esta ampliación permitirá analizar
	si el rendimiento depende del orden y de la selección concreta de los ejemplos
	proporcionados al modelo.
	
	\item \textbf{Construir un sistema experimental reproducible} que automatice la
	ejecución de los experimentos, el registro de las respuestas del modelo y el cálculo
	de métricas, permitiendo auditar cualquier resultado de forma independiente.
	
	\item \textbf{Evaluar las estrategias} sobre puzzles reales del benchmark
	\textsc{Linguini} en lenguas de diversas familias lingüísticas, en ambas direcciones
	de traducción (lengua desconocida $\rightarrow$ inglés e inglés $\rightarrow$ lengua
	desconocida), utilizando modelos de distintas familias y tamaños.
	
	\item \textbf{Comparar los resultados obtenidos} con los valores de referencia
	publicados en \citet{sanchez2024linguini} mediante las métricas BLEU, chrF++ y TER,
	e identificar en qué condiciones cada estrategia mejora o empeora respecto al
	baseline directo.
	
	\item \textbf{Realizar un análisis cualitativo} comparando el proceso de resolución
	humana de los puzzles con el razonamiento explícito generado por los modelos,
	identificando los puntos de divergencia y los patrones de error más frecuentes según
	la familia lingüística, el tamaño del modelo y la dirección de traducción.
	
\end{enumerate}

% ─────────────────────────────────────────────────────────────
\section{Restricciones y condicionantes}
\label{sec:restricciones}
% ─────────────────────────────────────────────────────────────

El desarrollo del trabajo está sujeto a las siguientes restricciones:

\begin{description}
	
	\item[Acceso a los modelos:] se utilizan modelos accesibles mediante APIs externas
	gratuitas, como Groq y OpenRouter. Esta decisión limita los modelos disponibles,
	pero reduce el coste computacional y facilita la reproducción de los experimentos.
	
	\item[Límites de tasa:] las APIs gratuitas imponen límites de peticiones por minuto y
	por día, lo que condiciona el número de experimentos que pueden ejecutarse en una
	sesión y obliga a introducir retardos entre llamadas. Este condicionante es
	especialmente notable en M2, cuya naturaleza iterativa genera varias llamadas
	consecutivas por puzzle, y será todavía más relevante en M2.2 debido al número de
	combinaciones posibles.
	
	\item[Coste combinatorio de M2.2:] la extensión combinatoria incrementa
	considerablemente el número de llamadas necesarias. Para un puzzle con $n$ pares de
	entrenamiento, el número de subconjuntos evaluados es:
	
	\[
	\sum_{k=2}^{n} \binom{n}{k}
	\]
	
	Por ejemplo, para $n=10$ serían necesarias 1013 iteraciones por puzzle y modelo. Este
	coste obliga a planificar cuidadosamente la selección de puzzles y proveedores.
	
	\item[Datos:] los puzzles utilizados provienen exclusivamente del benchmark
	\textsc{Linguini} \citep{sanchez2024linguini}, cuya descarga requiere aceptar los
	términos de uso del corpus de la IOL (licencia CC-BY-SA~4.0). No se utilizan puzzles
	de elaboración propia para mantener la comparabilidad con la literatura.
	
	\item[Datos de preentrenamiento desconocidos:] no es posible conocer con precisión
	todos los textos utilizados durante el entrenamiento de los modelos evaluados. Por
	este motivo, no puede garantizarse que las lenguas analizadas sean completamente
	desconocidas para el modelo. El control M0 se utiliza únicamente para detectar
	posibles indicios de conocimiento previo.
	
	\item[Ausencia de reentrenamiento:] no se realiza ningún tipo de \textit{fine-tuning}
	ni adaptación de los modelos. Todas las estrategias operan mediante prompting, sin
	modificar los pesos del modelo.
	
	\item[Variabilidad de las respuestas:] todos los experimentos se ejecutan con
	temperatura cero (\texttt{temperature=0}) para reducir la variabilidad entre
	ejecuciones. No obstante, no puede garantizarse un determinismo absoluto, ya que los
	proveedores externos pueden modificar internamente los modelos o su infraestructura
	de inferencia.
	
\end{description}